\section{Conclusions}\label{sec:con}
We introduced counterfactual querying as a new data management primitive: a constrained, cost-aware, and set-oriented retrieval problem over observed data. Guided by principles of interpretability (P1), cost faithfulness (P2), efficiency (P3), and collective validity (P4), we developed a two-phase framework that jointly considers these principles, highlighting the gap from similarity-based search. Our analysis across modalities and the \textbf{DcCA} prototype show this vision is feasible and scalable, positioning data systems as platforms for structured, actionable reasoning rather than simple retrieval. Elevating counterfactual querying to a first-class abstraction enables integrating constraints, costs, and decision objectives directly into query processing. Looking ahead, data-centric counterfactual analysis can support human-centered AI with interpretable and reliable systems. Realizing this requires advances in constraint-aware indexing, cost modeling, and set-level optimization.